\documentclass[11pt,a4paper]{article}

\usepackage[utf8]{inputenc}
\usepackage[spanish]{babel}
\usepackage{geometry}
\usepackage{hyperref}
\usepackage{enumitem}
\usepackage{titlesec}
\usepackage{booktabs}
\usepackage{array}
\usepackage{xcolor}
\usepackage{fancyhdr}
\usepackage{parskip}

\geometry{margin=2.5cm}

\definecolor{d}{RGB}{30, 58, 95}
\definecolor{a}{RGB}{185, 107, 60}
\definecolor{g}{RGB}{120, 120, 120}

\pagestyle{fancy}
\fancyhf{}
\rhead{\textcolor{g}{\small MexGo — docs internos}}
\lhead{\textcolor{g}{\small ARQUITECTURA.tex}}
\rfoot{\textcolor{g}{\small \thepage}}
\renewcommand{\headrulewidth}{0.3pt}

\titleformat{\section}{\normalsize\bfseries\color{d}}{}{0em}{}[\titlerule]
\titleformat{\subsection}{\normalsize\bfseries\color{g}}{}{0em}{}

\newcommand{\ok}[1]{\textcolor{teal}{\textbf{✓ #1}}}
\newcommand{\pend}[1]{\textcolor{red}{\textbf{? #1}}}
\newcommand{\dec}[1]{\textcolor{a}{\textbf{→ #1}}}
\newcommand{\descartado}{\textcolor{red}{\textbf{✗ Descartado}}}
\newcommand{\parcial}{\textcolor{a}{\textbf{~ Principio adoptado}}}
\newcommand{\elegido}{\textcolor{teal}{\textbf{✓ Elegido}}}
\newcommand{\f}[1]{\texttt{\textcolor{d}{#1}}}

\hypersetup{colorlinks=true, linkcolor=d, urlcolor=a,
  pdftitle={MexGo ARQUITECTURA}, pdfauthor={Los Mossitos}}

% ===========================================================================
\begin{document}
% ===========================================================================

{\centering
  {\Large\bfseries\color{d} MexGo — Arquitectura}\par
  \vspace{0.3cm}
  {\small\color{g} Los Mossitos · Genius Arena 2026 · \today}\par
  \vspace{0.5cm}
  \hrule
  \vspace{0.8cm}
}

\tableofcontents
\vspace{1cm}
\hrule
\vspace{0.5cm}

% ===========================================================================
\section{Para qué sirve este doc}
% ===========================================================================

Decisiones de arquitectura del proyecto. Qué se evaluó, qué se eligió y por
qué. Actualizar cuando cambie algo. Responsable: Héctor.

% ===========================================================================
\section{Patrones evaluados}
% ===========================================================================

\subsection{Monolito modular — \descartado\ como patrón puro}

Todo en un repo Next.js. Rápido de arrancar y fácil de desplegar, pero con
5 personas sin separación clara los conflictos de git matan el hackathon.
El algoritmo de equidad terminaría mezclado con React si no se separa a
propósito. \dec{Se adopta solo la ventaja del repo único.}

\subsection{Cliente / API separados — \descartado\ como implementación}

Frontend y backend en proyectos distintos. Separación limpia, pero en 2 días
el setup de CORS, dos repos y contratos de API come tiempo real. Si el backend
no levanta, el frontend no prueba nada. \dec{Se adopta como mentalidad
de separación, no como implementación.}

\subsection{BFF — Backend for Frontend — \elegido\ como mentalidad}

Un intermediario orquesta Gemini + Supabase + algoritmo y devuelve al
frontend el resultado listo. El frontend solo renderiza. El algoritmo nunca
llega al navegador. Requiere definir contratos de API el día 1.
\dec{Implementado con Next.js API Routes, no como servicio separado.}

\subsection{Serverless / Supabase Edge Functions — \descartado\ para hackathon}

Funciones Deno en la infra de Supabase. Escalan solas, sin servidor propio.
Deno es casi idéntico a Node.js pero las diferencias en imports y el cold
start de 1--2s añaden fricción innecesaria en 2 días. \dec{Ruta de migración
natural post-hackathon.}

\subsection{Microservicios — \descartado}

Un servicio por dominio, cada uno con su BD. Escalable y tolerante a fallos,
pero service discovery + load balancing + trazabilidad distribuida para 5
personas en 2 días es absurdo. \dec{No aplica.}

\subsection{Clean Architecture — \parcial}

La lógica de negocio no conoce a Supabase ni a Gemini. Elegante y testeable,
pero el boilerplate de interfaces y capas formales no cabe en 2 días.
\dec{El algoritmo y el perfil cultural son funciones puras en \f{lib/}
que no importan proveedores directamente.}

\subsection{Arquitectura hexagonal — \parcial}

Puertos y adaptadores. Cambiar Gemini por OpenAI no toca la lógica de negocio.
Misma ventaja que Clean, mismo problema de boilerplate.
\dec{Gemini, Mapbox y Cloudinary solo se acceden desde \f{lib/}.
Si cambia un proveedor, solo cambia ese archivo.}

% ===========================================================================
\section{Patrón elegido}
% ===========================================================================

\elegido\ \textbf{Monolito modular con mentalidad BFF — Next.js API Routes.}

Un repo. Las API Routes son el BFF: reciben peticiones del frontend,
llaman a Gemini y Supabase, aplican el algoritmo, devuelven tarjetas listas.
El frontend solo muestra. La lógica de negocio vive en \f{lib/} como
funciones puras en TypeScript, sin depender de Next.js ni de ningún proveedor.

Por qué este: Next.js tiene curva de aprendizaje factible, un repo elimina coordinación de
servicios, las API Routes protegen el algoritmo del cliente, y post-hackathon
se puede migrar a Edge Functions sin tocar el frontend.

\subsection{Reglas — no se rompen}

\begin{itemize}[leftmargin=*, itemsep=2pt]
  \item Frontend no importa Supabase directo — solo consume \f{/app/api/}.
  \item Algoritmo solo en \f{lib/equity.ts} — nunca en un componente React.
  \item Claves de Gemini, Mapbox server y Supabase service role — nunca en
        \texttt{NEXT\_PUBLIC\_}.
  \item Cada archivo en \f{lib/} — una sola responsabilidad.
\end{itemize}

% ===========================================================================
\section{Estructura de carpetas}
% ===========================================================================

\begin{verbatim}
mexgo/
├── app/
│   ├── (turista)/
│   │   ├── page.tsx              # Bienvenida
│   │   ├── cuestionario/
│   │   ├── itinerario/
│   │   └── mapa/
│   ├── (admin)/
│   │   ├── login/
│   │   ├── dashboard/
│   │   └── negocio/[id]/
│   └── api/                      # BFF — nunca llamado directo desde cliente
│       ├── recomendar/           # POST: perfil → tarjetas
│       ├── itinerario/           # POST: preferencias → itinerario
│       ├── negocios/             # CRUD negocios
│       └── auth/
│
├── lib/                          # Lógica pura — sin deps de Next.js
│   ├── equity.ts                 # Algoritmo de equidad
│   ├── cultural.ts               # Perfil cultural por país
│   ├── gemini.ts                 # Wrapper Gemini
│   ├── mapbox.ts                 # Wrapper Mapbox
│   ├── supabase.ts               # Cliente servidor
│   └── supabase-client.ts        # Cliente browser (solo auth)
│
├── components/
│   ├── ui/
│   ├── turista/
│   └── admin/
│
├── docs/
├── supabase/migrations/
├── .env.local                    # Nunca en git
├── .env.example                  # Sí en git — sin valores reales
└── next.config.ts
\end{verbatim}

% ===========================================================================
\section{Flujo principal}
% ===========================================================================

\begin{enumerate}[leftmargin=*, itemsep=2pt]
  \item Turista abre PWA por QR. Next.js detecta idioma con
        \texttt{Accept-Language}.
  \item Cuestionario → \texttt{POST /api/recomendar} con perfil básico.
  \item API Route → \f{lib/cultural.ts} (parámetros del país) →
        \f{lib/equity.ts} (score con negocios de Supabase) → 4--6 tarjetas.
  \item Frontend renderiza. No procesa nada.
  \item Si tiene boleto → \texttt{POST /api/itinerario} → Gemini integra
        el partido al plan del día.
  \item Itinerario se guarda en \texttt{localStorage} para offline.
\end{enumerate}

% ===========================================================================
\section{Responsables}
% ===========================================================================

\begin{center}
\begin{tabular}{>{\bfseries}p{2.5cm} p{2.5cm} p{8.5cm}}
\toprule
Módulo & Quién & Archivos \\
\midrule
Arquitectura + deploy & Héctor & \f{next.config.ts}, AWS, \f{.env} \\
Algoritmo             & Héctor & \f{lib/equity.ts}, \f{lib/cultural.ts} \\
Gemini                & Héctor & \f{lib/gemini.ts}, \f{api/itinerario/} \\
Frontend turista      & Emi    & \f{app/(turista)/}, \f{components/turista/} \\
Frontend admin        & Farid  & \f{app/(admin)/}, \f{components/admin/} \\
Backend / Supabase    & Alan   & \f{app/api/}, \f{lib/supabase.ts}, migraciones \\
Apoyo / QA            & Xavier & datos de prueba, tests \\
\bottomrule
\end{tabular}
\end{center}

% ===========================================================================
\section{Pendientes}
% ===========================================================================

\begin{itemize}[leftmargin=*, itemsep=2pt]
  \item \pend{Servicio AWS} — Amplify / EC2+Nginx / App Runner. (Héctor)
  \item \pend{Contratos de API} — schemas por endpoint antes del día 1.
        Alan + Héctor → \f{SCHEMA.tex}.
  \item \pend{Caché del score} — memoria / Supabase / Redis.
        Impacta \f{lib/equity.ts}. (Héctor)
\end{itemize}

% ===========================================================================
\section{Cambios}
% ===========================================================================

\begin{tabular}{p{2cm} p{2.5cm} p{9cm}}
\toprule
Fecha & Quién & Qué \\
\midrule
\today & Héctor & v1.0 — evaluación de patrones, decisión inicial. \\
\today & Héctor & v1.1 — nombres reales, párrafos cortos. \\
\today & Héctor & v1.2 — estilo cuaderno, sin relleno. \\
\bottomrule
\end{tabular}

\end{document}
